\section{Introduction}\label{sec:intro}

Neurosymbolic systems combine neural perception with symbolic reasoning, but established integrations often place the two halves on different processors. Deep learning frameworks---PyTorch, JAX---execute dense tensor computation on the GPU through optimized CUDA kernels, while systems such as DeepProbLog~\cite{deepproblog} and NeurASP~\cite{neurasp} use CPU-side symbolic backends. Their integration can therefore cross the PCIe bus for network outputs, symbolic results, or gradients. At scale, repeated host--device movement can dominate wall-clock time and bandwidth.

The transfer wall is one half of the problem; fragmented interfaces are the other. Existing neurosymbolic frameworks commonly bridge a neural network to one symbolic backend (DeepProbLog to probabilistic logic, NeurASP to answer-set programming). GPU logic efforts, conversely, accelerate focused workloads such as deterministic Datalog~\cite{gpulog,vflog} or SAT~\cite{parafrost}. xlog explores how a typed language and shared CUDA services can support deterministic, probabilistic, epistemic, and neural-symbolic use without pretending that those routes have identical execution or differentiability properties.

xlog is a CUDA-native logic programming language whose compiler and runtime treat a neural network as an ordinary predicate. Network outputs can become weighted facts for probabilistic knowledge compilation, where exact query gradients flow back to the network. Other reasoning modes reuse the typed relational and solver interfaces without claiming the same gradient semantics. Device-resident buffers carry facts, deltas, circuits, solver state, and gradients, while the host performs orchestration, I/O, compilation, and bounded control reads. The certified resident recursive and Monte Carlo sampled cores provide the stronger measured boundary: no tracked host--device transfers inside the core, followed by one terminal receipt. Verified knowledge compilation certifies the final smoothed circuit against its back-end formula before use, and DLPack~\cite{dlpack} exports device gradient buffers directly into PyTorch's autograd graph.

The principal contributions of this paper are:

\begin{itemize}
  \item \textbf{A shared neurosymbolic integration model.} A neural network is a predicate, and typed device-buffer interfaces connect it to the supported reasoning routes. End-to-end gradient training is realized specifically through the probabilistic knowledge-compilation path (\Cref{sec:arch,sec:nesy}).
  \item \textbf{CUDA relational execution with a certified resident route.} Semi-naive Datalog evaluation with stratified negation and aggregation uses custom CUDA relational kernels. Eligible recursive query plans are selected automatically for a resident conditional graph whose measured core has no tracked transfers; the wider ordinary executor remains host-orchestrated. A worst-case-optimal join (WCOJ) subsystem avoids the intermediate-relation blowup of binary-join plans on skewed cyclic queries (a measured $27.96\times$ geometric-mean ablation, \Cref{sec:datalog}).
  \item \textbf{Verified knowledge compilation.} A CUDA-backed provenance $\to$ CNF $\to$ Decision-DNNF $\to$ circuit pipeline certifies the final smoothed circuit against its back-end source formula before caching or evaluation. The host observes bounded solver status and error scalars; resolution-proof traces are checked on the GPU. Structural invariants required for weighted model counting are established separately by construction (\Cref{sec:prob}).
  \item \textbf{End-to-end differentiable training with zero-copy interop.} Compiled circuits are cached across iterations and evaluated as level-parallel forward/backward kernels, with gradients exported zero-copy via DLPack and Arrow; circuit caching yields a measured $2.74\times$ training speedup. Term embeddings and differentiable ILP ride the same device-resident path (\Cref{sec:nesy}).
  \item \textbf{Perception learned through symbolic credit on an external benchmark.} On the CAVIAR video-event corpus, under a direct per-timestep target, the rule search returns the same theory whether its proximity predicate is the hand-set geometric one that dedicated rule learners are handed or a network trained by the logic's own credit alone, never shown a distance label or any target of its own, at no cost in held-out accuracy. Event-Calculus rules over the precomputed predicate are induced under pre-registered gates; substituting the learned detector inside \emph{that} search is where the result currently fails (\Cref{sec:ecinduction}).
  \item \textbf{Rule induction at scale.} On the Brest AIS maritime corpus (3{,}548 gold \texttt{rendezVous} intervals) the same protocol, over pair-level folds and permutation-null gates, measures in isolation the two mechanisms the published Event-Calculus learners are built on. Weighting the clauses of the same gated pool raises micro F1 by $0.065$ over crisp selection, against our own prior expectation that it would not, and learning those weights in a single chronological pass reproduces the batch predictions exactly, for a reason we can name. We draw no comparison with the published maritime figure, whose temporal grid differs from ours (\Cref{sec:maritime}).
  \item \textbf{Reasoning beyond probability.} Epistemic reasoning over world views (\texttt{know}/\texttt{possible} under founded and Gelfond-style semantics) and well-founded semantics execute on the same substrate over finite supported programs, including positive modal fixpoints and supported recursion through negation evaluated by GPU-backed well-founded semantics, reusing the CDCL, join, and circuit machinery instead of a separate engine (\Cref{sec:epistemic}). The accepted fragment and its bounds are in \Cref{sec:limits}.
\end{itemize}

\textbf{Target workloads.} Avoiding repeated data-plane copies matters most where the symbolic workload is large and repeated: program-analysis queries over millions of ground atoms, probabilistic logic with many training iterations that share one compiled circuit, and knowledge-graph reasoning with trainable entity embeddings. Small tasks such as MNIST addition exercise the integration end-to-end but sit below this regime. They demonstrate correctness and the circuit-caching effect, with the measured transfer component secondary: approximately 10\% of a training step at a standard batch and lost in run-to-run variation when a step makes only a handful of neural$\to$symbolic handoffs (\Cref{sec:overhead}). Accordingly, xlog targets NVIDIA-GPU workloads whose byte working set and fixed capacities fit the selected route (\Cref{sec:limits}) and that need neural perception coupled to exact, verified probabilistic inference rather than fuzzy relaxations.

\Cref{sec:arch} presents the integration model and system architecture. \Cref{sec:lang} describes the xlog language. \Cref{sec:datalog} details the GPU-resident symbolic substrate. \Cref{sec:prob} presents verified knowledge compilation, the central technical result. \Cref{sec:nesy} ties the pieces together as end-to-end neurosymbolic learning. \Cref{sec:ecinduction} puts that surface on an external video-event benchmark, and \Cref{sec:maritime} carries the same induction protocol to a larger maritime corpus. \Cref{sec:epistemic} extends the substrate to epistemic and other reasoning modes. \Cref{sec:eval} reports evaluation results, \Cref{sec:related} surveys related work, \Cref{sec:limits} discusses limitations and future work, and \Cref{sec:conclusion} concludes.
